% CS 418: Intro to Data Science — Worksheet 8 (08-week04-day2)
% Week 4, Wednesday 2026-09-16
% Compile with: latexmk -pdf worksheet.tex

\documentclass[11pt]{article}

% ── Packages ──────────────────────────────────────────────────────────────────
\usepackage[margin=0.65in, headheight=14pt]{geometry}
\usepackage{amsmath, amssymb}
\usepackage{ragged2e}
\usepackage{enumitem}
\usepackage{booktabs}
\usepackage{graphicx}
\usepackage{hyperref}
\usepackage{xcolor}
\usepackage{fancyhdr}
\usepackage{mdframed}
\usepackage{array}

% ── Page style ────────────────────────────────────────────────────────────────
\pagestyle{fancy}
\fancyhf{}
\lhead{\textbf{CS 418: Intro to Data Science}}
\rhead{UIC, Fall 2026}
\cfoot{\thepage}
\renewcommand{\headrulewidth}{1pt}
\setlength{\headheight}{12pt}
\setlength{\headsep}{0.4cm}

% ── Hyperlinks ────────────────────────────────────────────────────────────────
\hypersetup{colorlinks=false, hidelinks}

% ── Convenience environments ──────────────────────────────────────────────────
% Usage: \blankspace{6cm} — blank area for hand-written work, no rules
\newcommand{\blankspace}[1]{%
  \vspace{0.3em}
  \begin{minipage}[t][#1][t]{\linewidth}
  \end{minipage}
  \vspace{0.3em}
}

% Usage: \guidedopen{title}{guided content}{guided workspace height}{open-ended workspace height}
% Left: a titled, structured prompt with its own blank workspace. Right: a
% fully blank workspace for open-ended exploration of the same topic (no
% prompt text), sized so the two columns land at roughly the same total
% height. Neither workspace is ruled.
\newcommand{\guidedopen}[4]{%
  \noindent\textbf{#1}\\[0.18em]
  \noindent
  \begin{minipage}[t]{0.485\textwidth}
    #2
    \blankspace{#3}
  \end{minipage}
  \hfill
  \begin{minipage}[t]{0.485\textwidth}
    \blankspace{#4}
  \end{minipage}
  \\[0.4em]
}

% Usage: \panelbox{width}{height}{caption} — a framed drawing panel with `caption`
% centered beneath it. An empty caption leaves the panel bare, with nothing
% below the frame at all.
\newcommand{\panelbox}[3]{%
  \begin{minipage}[t]{\dimexpr#1+2\fboxsep+2\fboxrule\relax}
    \framebox[#1][c]{\rule{0pt}{#2}}%
    \if\relax\detokenize{#3}\relax\else
      \\[0.25em]{\centering\small #3\par}%
    \fi
  \end{minipage}%
}

% ══════════════════════════════════════════════════════════════════════════════
\begin{document}
\RaggedRight

% ── Title block: topic left, info box right ───────────────────────────────────
\noindent
\begin{minipage}[t]{0.55\textwidth}
  \vspace{0pt}
  {\LARGE \textbf{Worksheet \#\,8}} \\[0.18em]
  {\large \textit{Week 4, Day 2: Scraping guidance, Bayes' law, more descriptive statistics, why visualize}}
\end{minipage}
\hfill
\begin{minipage}[t]{0.40\textwidth}
  \vspace{0pt}
  \small
  Name: \underline{\hspace{4.5cm}} \\[0.35em]
  NetID: \underline{\hspace{4.3cm}} \\[0.35em]
  Date: \underline{\hspace{4.4cm}} \\
\end{minipage}

\vspace{0.6em}
\noindent\rule{\linewidth}{0.6pt}
\vspace{0.1em}

\noindent\underline{\textbf{Guided checkpoints (optional)}}
\vspace{0.8em}

% ── 1. Six steps to scrape ───────────────────────────────────────────────────
\guidedopen{1. Six Steps to Scrape a Website}{
List the six typical steps of scraping a web page.
}{5.0cm}{5.25cm}

\vfill

% ── 2. Anatomy of Bayes' theorem ─────────────────────────────────────────────
\guidedopen{2. Bayes' Anatomy}{
Write Bayes' law, define prior and posterior
}{4.6cm}{5.3cm}

\vfill

% ── 3. The EDA toolkit ───────────────────────────────────────────────────────
\guidedopen{3. Review: The EDA Toolkit}{
Name the methods (and attribute) that offer an exploratory look at a dataframe.
}{4.35cm}{5.55cm}

\vfill

\newpage

% ── 4. The shape of a distribution ───────────────────────────────────────────
\guidedopen{4. Common Shapes of Distributions}{
Sketch a distribution in each panel, mark the mean and median, and label the panel.

\vspace{0.5em}

\noindent
\panelbox{2.5cm}{2.5cm}{}\hfill
\panelbox{2.5cm}{2.5cm}{}\hfill
\panelbox{2.5cm}{2.5cm}{}
}{0.2cm}{5.05cm}

\vfill

% ── 5. Chebyshev's bounds ────────────────────────────────────────────────────
\guidedopen{5. Chebyshev's Bounds, Tabulated}{
Fill in the bound for each $k$.

\vspace{0.4em}

{\small\renewcommand{\arraystretch}{1.75}%
\noindent\begin{tabular}{@{}c >{\centering\arraybackslash}p{2.9cm} >{\centering\arraybackslash}p{2.9cm}@{}}
\toprule
$k$ & \textbf{Within $k$ SDs} & \textbf{Outside $k$ SDs} \\
\midrule
1 & & \\
2 & & \\
3 & & \\
4 & & \\
5 & & \\
\bottomrule
\end{tabular}}
}{0.4cm}{6.27cm}

\vfill

% ── 6. Anscombe's quartet ────────────────────────────────────────────────────
\guidedopen{6. Anscombe's Quartet}{
Sketch each of the four datasets.

\vspace{0.5em}

\noindent
\panelbox{4.0cm}{2.3cm}{I}\hfill
\panelbox{4.0cm}{2.3cm}{II}

\vspace{0.4em}

\noindent
\panelbox{4.0cm}{2.3cm}{III}\hfill
\panelbox{4.0cm}{2.3cm}{IV}

\vspace{1.2em}

\noindent All four have...
}{1.3cm}{9.2cm}

\vfill

\end{document}
